\section{Generated Code}
\label{sec:generatedcode}

The generators create one versioned Deliverable Set for the estate, laid out
per project and application by the path plan the resolved model carries.
Table~\ref{tab:generated_resources} lists its resources and how they are
derived.

\begin{table}[!htbp]
\centering
\small
\begin{tabular}{>{\raggedright\arraybackslash}p{0.24\linewidth}>{\raggedright\arraybackslash}p{0.68\linewidth}}
\hline
\textbf{Output} & \textbf{How it is derived} \\
\hline
\texttt{Namespace} and indexes &
  A named cluster partition for the project and index files listing what to deploy. \\
Workload resources &
  Definitions that run processes, give each an identity and keep enough instances available during maintenance. \\
\texttt{NetworkPolicy} &
  Firewall rules that allow required process connections and block all others. \\
\texttt{IngressRoute} &
  Web proxy rules that send each address and path to the correct process. \\
Monitoring and alerts &
  Instructions for collecting metrics and conditions for raising alerts. \\
Secret resources &
  Vault permissions granting each process only its allowed secrets. \\
\hline
\end{tabular}
\caption{Generated Resources and Their Derivation}
\label{tab:generated_resources}
\end{table}

The transformation computes limits, startup probes, route precedence and policy
peers. For \texttt{auth}, 96 declared lines produce eleven files of about 620
lines. The repository contains hand-written expected files, which are used to
check the generated result. The tree carries one index entry point at its root,
so the whole estate is applied through a single reference rather than one per
application.
